% Results section, draft of 5 Nov 2026 (Claude). Every number comes from a saved run; the folder is in the comment
% next to it. Validation = the 60 rooms of scenarios_rooms_val.json, test = the 80 rooms of scenarios_rooms_test.json
% (run once, on 5 Nov, after every choice was made). Needs \usepackage{booktabs} or swap \toprule etc. for \hline.

\section{Results}
\label{sec:results}

\subsection{Metrics}
The base is scored as \emph{parked} when, for the whole of the last \SI{2}{s} of the episode, it stands where the
arm can reach the target (horizontal distance from \SI{0.3}{m} up to 84.9\% of the straight arm's reach when facing
the target, falling linearly to 83\% at the \SI{30}{\degree} facing limit; both edges were measured with the expert
arm), faces the target within \SI{30}{\degree}, is stopped (below \SI{0.02}{m/s} and \SI{0.05}{rad/s}) and has had no
contact in the whole run. A mover that hits a parked robot also counts as a failure. We report mover hits separately,
because they were the main failure throughout. The full robot uses the course rule: the hand stays within \SI{1}{cm}
of the target for \SI{0.25}{s}, with no contact in the whole run. Base runs last \SI{70}{s}; full-robot runs last
\SI{100}{s}, which allows for the unfolding and the reach.

\subsection{Learning the room base}
Table~\ref{tab:progression} follows the room base through the changes that were kept. Each row is the third DAgger
round~\cite{ross2011dagger} unless marked otherwise. With the visible-only RRT teacher, removing the velocity inputs
raised parking from 4 to 9 rooms; with them the network copied its own motion (the copycat problem~\cite{wen2020copycat}).
A small learned 1-D convolution over the lidar rays (circular, shared across directions) raised parking to 16, and the threat head (an auxiliary output predicting
the nearest mover's position and velocity) lowered mover hits to 30. The largest gains came from changing what the
network learns from: the dynamic-window teacher, which the network could copy far better than the RRT (copy $R^2$ 0.46
against 0.06 on the same rooms), and folding the arm inside the base outline while driving. The old arm pose reached
\SI{0.73}{m} ahead below the obstacle tops, where movers struck it. DAgger rounds that added crossing-mover scenes cut
mover hits from 16 to 12. A small learned ``hold still'' head on the frozen network, latched once it has been on for
\SI{0.5}{s}, stops the slow creeping that failed the stillness check. The final room base has 10.9k parameters.

\begin{table}[t]
\centering
\caption{Room base on the 60 validation rooms (parked; runs with a mover contact).}
\label{tab:progression}
\begin{tabular}{lcc}
\toprule
Model (DAgger round 3 unless noted) & Parked & Mover hits \\
\midrule
Visible-only RRT teacher, velocity inputs, 1 scan         & 4  & 45 \\ % dagger-rooms-velocity/r3
\quad no velocity inputs                           & 9  & 37 \\ % dagger-rooms-1scan/r3
\quad 3 scans                                      & 7  & 36 \\ % dagger-rooms-3scans/r3
\quad + ray convolution                            & 16 & 35 \\ % dagger-rooms-conv/r3
\quad + threat head                                & 15 & 30 \\ % dagger-rooms-conv-threat/r3
Window teacher, 72 rays                            & 21 & 23 \\ % dagger-rooms-window72b/r3
\quad arm folded while driving (no retraining)     & 31 & 16 \\ % window-fold-r0-start/validation
\quad + crossing-mover DAgger                      & 35 & 12 \\ % dagger-rooms-crossing/r3
\quad + latched hold head (\textbf{final})         & \textbf{37} & \textbf{12} \\ % crossing-hold/r3-hold-latch-validation
\midrule
Second seed: crossing DAgger / + hold              & 32 / 31 & 13 / 13 \\ % dagger-rooms-crossing-s1/r3, crossing-hold/s1-r3-hold-latch-validation
600 more rooms + mirror augmentation, 2 seeds      & 28, 36 & 15, 15 \\ % dagger-rooms-big-s0/-s1 r3
Teacher (window expert)                            & 50 & 4 \\ % rooms-window72-val
\bottomrule
\end{tabular}
\end{table}

\textbf{Run-to-run noise.} Two seeds of the same crossing DAgger disagree on 19 of the 60 rooms, so differences of
2--4 parked rooms are within noise. The drop in mover hits (16 to 12--13) and the better copying just before past
mover hits (turn correlation 0.21 to 0.53--0.55) repeated on both seeds. The parking gain from the crossing DAgger and
the hold head did not repeat clearly; the hold head stayed in the final model because the rule we wrote before the run
selected it. Several ideas were tried and not kept because they failed rules written before the results: a mixture
output and a menu of discrete actions (no better copying in the states before hits), a learned evade/target switch,
a larger network (copying on new rooms did not improve: 0.65 / 0.70 / 0.66 for 64 / 128 / 256 units), a fourth DAgger
round (35 parked, 14 mover hits), and four times more training rooms with left--right mirror augmentation (mover
hits rose to 15 on both seeds).

\subsection{Final test}
The test rooms were used once, after the base, the arm and every evaluation setting were fixed
(Table~\ref{tab:test}). Both results agree with validation within noise, so the choices made on validation did not
overfit noticeably.

\begin{table}[t]
\centering
\caption{Final models on the validation and test rooms.}
\label{tab:test}
\begin{tabular}{lcc}
\toprule
 & Validation (60) & Test (80) \\
\midrule
Base parked                            & 37 (62\%) & 48 (60\%) \\ % crossing-hold/r3-hold-latch-validation; final-test/base
Base: runs with a mover hit            & 12 (20\%) & 18 (23\%) \\
Full robot success                     & 32 (53\%) & 45 (56\%) \\ % arm-eval-rooms-full-val; arm-eval-final-test
Arm success after a clean hand-off     & 32/38 (84\%) & 45/60 (75\%) \\
\bottomrule
\end{tabular}
\end{table}

For the full robot, the arm starts only once the base passes the parking rule for \SI{2}{s}. The folded arm first
moves along a straight joint-space path to its usual start pose (about \SI{7}{s}); the learned arm then reaches in a
median of \SI{2.8}{s}. Of the 35 test failures, 16 were mover contacts while driving (10 of these robots still parked
afterwards), 4 never passed the parking rule and 1 parked too late to reach. The other 14 were arm misses without
contact: in 12 the hand came within 1.0--\SI{3.5}{cm} but did not stay within \SI{1}{cm}. The same two failure types
dominate validation, where unfolding caused a single contact.

\subsection{Open floor and compute}
On the open-floor task with static obstacles (40 held-out scenarios), the learned robot (a 7k-parameter base
and a 6.2k-parameter arm) succeeds in 37 of 40 with no contact; the expert succeeds in 40 of 40, as does expert
driving with the learned arm. % arm-eval-final, -hybrid, -expert
Both networks run on one CPU thread (AMD Ryzen 7 5800X3D): the room base takes \SI{178}{\micro s} per step
and the arm \SI{49}{\micro s}, well inside the \SI{1}{ms} control period. % scratchpad efficiency check, 5 Nov

\subsection{Limitations}
The teacher is privileged: it knows the map of walls and still obstacles, and for every mover a lidar ray hit in the
last second it uses the mover's true position and velocity. The network sees only lidar and its own state, so part of
the gap to the teacher (50 against 37 parked) is information the network does not have. Most comparisons use
a single seed, and the noise above limits what they can show. Mover contacts while driving remain the main loss,
followed by the arm's last centimetre.
